% !PRINTED 85
\weberpage{85}
%<<CONT
der Faktor $Q$ bestimmen. Man erh\"alt zun\"achst
\[
1 = Q^2\prod^{\nu}\frac{(1-q^{2\nu-1})^2(1+q^{2\nu-1})^2(1+q^{2\nu})^2}{(1-q^{2\nu})^2},
\]
oder, indem man \"uber das noch unbestimmte Vorzeichen von $Q$ und damit \"uber die
\gesperrt{Vorzeichen} der $\vartheta$-Funktionen verf\"ugt:
\[
Q = \prod_{1,\infty}^{\nu}\frac{(1-q^{2\nu})}{(1+q^{2\nu})(1+q^{2\nu-1})(1-q^{2\nu-1})}.
\]
Im Nenner kann man f\"ur $\prod(1+q^{2\nu})(1+q^{2\nu-1})$ setzen $\prod(1+q^{\nu})$, und
der Z\"ahler $\prod(1-q^{2\nu})$ l\"a\ss{}t sich zerlegen in
\[
\prod(1-q^{\nu})(1+q^{\nu}) = \prod(1-q^{2\nu})(1-q^{2\nu-1})(1+q^{\nu}).
\]
Dadurch ergibt sich endlich
\weq{6}{Q = \prod^{\nu}(1 - q^{2\nu}).}
Die Ausdr\"ucke (1), (2), (3), (4) lassen sich in reeller Form darstellen, wenn man die
Multiplikation der konjugiert imagin\"aren Faktoren ausf\"uhrt. Man erh\"alt so:
\weq{7}{\begin{gathered}
\vartheta_{00}(v) = \prod_{1,\infty}^{\nu}(1 - q^{2\nu})(1 + 2q^{2\nu-1}\cos 2\pi v + q^{4\nu-2}),\\
\vartheta_{01}(v) = \prod_{1,\infty}^{\nu}(1 - q^{2\nu})(1 - 2q^{2\nu-1}\cos 2\pi v + q^{4\nu-1}),\\
% [sic] Printed exponent 4\nu-1; mathematically expected 4\nu-2.
\vartheta_{10}(v) = 2q^{\frac{1}{4}}\cos\pi v\prod_{1,\infty}^{\nu}(1 - q^{2\nu})(1 + 2q^{2\nu}\cos 2\pi v + q^{4\nu}),\\
\vartheta_{11}(v) = 2q^{\frac{1}{4}}\sin\pi v\prod_{1,\infty}^{\nu}(1 - q^{2\nu})(1 - 2q^{2\nu}\cos 2\pi v + q^{4\nu}).
\end{gathered}}
F\"uhren wir den Ausdruck (6) in die letzte Gleichung (5) ein, so ergibt sich
\[
\vartheta_{11}' = 2\pi q^{\frac{1}{4}}\prod(1 - q^{2\nu})^3,
\]
und wenn wir also
\weq{8}{\eta(\omega) = q^{\frac{1}{12}}\prod(1 - q^{2\nu})}
setzen, so wird
\weq{9}{\vartheta_{11}' = 2\pi\eta(\omega)^3.}
Indem wir $Q$ mit Hilfe der Relation
\[
Q = q^{-\frac{1}{12}}\eta(\omega)
\]
aus (5) eliminieren, setzen wir
\weq{10}{\begin{gathered}
\vartheta_{00} = \eta(\omega)f(\omega)^2,\\
\vartheta_{01} = \eta(\omega)f_1(\omega)^2,\\
\vartheta_{10} = \eta(\omega)f_2(\omega)^2,
\end{gathered}}
%CONT>>
